
%%--------------------------------------------------------------------------
%%--------------------------------------------------------------------------
\subsection{Examen de IST-SARO, 24 de mayo de 2024}

Se quiere construir un sitio web, MemeMeMola, donde puedes mandar imágenes de memes. Estas imágenes se pueden comentar y votar del 1 al 5. La funcionalidad básica del sitio es la siguiente:

\begin{enumerate}
  \item En el sitio no hay cuentas para usuarios: toda la funcionalidad está disponible para cualquiera que lo visite, salvo para subir una imagen de un meme nueva para lo que se necesitará un nombre de usuario.
  \item La página principal del sitio tiene un formulario que permitirá a los visitantes escoger un nombre de usuario. Puede haber varios visitantes con el mismo nombre. Una vez elegido el nombre, el nombre aparecerá en todas las paǵinas del sitio con el texto ``eres \emph{nombre} en MemeMeMola'', y debajo habrá un botón para \emph{salir} (esto es, dejar de usar ese nombre).
  \item La página principal del sitio mostrará un listado con las imágenes de los diez últimos memes que se han subido. Este listado incluirá, para cada meme: la imagen (como enlace a la página del meme, ver más abajo), el nombre de usuario que lo subió, el número de votos, la media de los votos (un valor numérico entre 1 y 5) y el número de comentarios que ha recibido ese meme. 
  \item Después del listado en la página principal, y si el visitante ha elegido nombre, aparecerá un formulario donde podrá subir un meme, introduciendo en un único campo el archivo con la imagen.
  \item Al subir la imagen de un meme, el visitante verá la página del meme que acaba de crear (por lo que cada meme tendrá su propia página). En esa página aparecerá la imagen del meme, su fecha de subida, el número de votos que ha recibido, la media de los votos recibidos y un formulario para poner un comentario. Los comentarios aparecerán en orden inverso de publicación (los más recientes primero).
  \item Cada visitante puede poner un comentario (habiendo elegido nombre o no) utilizando el formulario indicado anteriormente. Al poner un comentario, el visitante volverá a ver la página del meme que estaba viendo, pero ahora con el comentario que acaba de poner. Si no había elegido nombre, el comentario aparecerá como anónimo.
  \item Tras votar un meme, el visitante volverá a ver la página del meme que acaba de votar, ya con su voto considerado en la cuenta y en la media de votos. No se permitirá dos votos del mismo visitante.
  \item Las imágenes de los memes están alojados como un fichero en el propio sitio web de MemeMeMola.
  \item Todas las páginas del sitio tendrán una imagen de cabecera (banner), que será servida por el propio sitio.
  \item Todas las páginas del sitio tendrán una imagen de un píxel servido por el servicio \emph{itrackyou.com}.
\end{enumerate}

\newpage

Teniendo en cuenta los requisitos anteriores, se pide:

\begin{enumerate}
  \item Diseña un esquema REST para proporcionar el servicio descrito. Se habrán de especificar los nombres de recurso empleados, y cómo reaccionará la aplicación cuando reciba los métodos POST o GET sobre esas URLs (no se usarán los métodos PUT o DELETE). \textbf{Muy importante:} Coloca la información en una \textbf{tabla}, con los nombres de los recursos en una columna, los métodos en otra, la descripción de lo que realizará la aplicación al recibirlos en la tercera, y el contenido de los datos de la petición, si los hay en la cuarta (se consideran datos de la petición a los que vayan, normalmente como \emph{query string}, en el cuerpo de la petición HTTP o tras el nombre de recurso en la primera línea de la cabecera). Escribe también un fichero similar al fichero urls.py de Django (aunque no es importante que se respete la sintaxis mientras se entienda y la estructura sea similar a la de Django), que refleje el esquema REST anterior. (1 punto)

  \item Describe el modelo de datos que necesitará esta aplicación. Define las tablas necesarias y los campos necesarios para la funcionalidad descrita. Hazlo de forma lo más similar posible a lo que tendrías que escribir en el fichero models.py en Django (aunque no es importante que se respete la sintaxis mientras se entienda el modelo de datos que propones). (1 punto)

  \item Describe las interacciones HTTP que ocurrirán entre el navegador y cualquier servidor web en el siguiente escenario: El escenario comienza cuando un visitante que accede por primera vez al sitio (pone en el navegador la URL de la página principal del sitio). A continuación, después de ver esta página principal, elige un nombre y sube un nuevo meme. El escenario termina cuando el visitante ve la página del meme que acaba de crear. (1 punto)

  \item Describe las interacciones HTTP que ocurrirán entre el navegador y cualquier servidor web en el siguiente escenario: El escenario comienza cuando un visitante que accede por primera vez al sitio (pone en el navegador la página de un meme en concreto). A continuación, después de ver esta página, le da al meme un voto de un 5. Finalmente, añade un comentario que dice ``buenísimo XD''. El escenario termina cuando el visitante ve la página del meme con el voto y el comentario. (1 punto)

  \item MemeMeMola tiene mucho éxito y te gustaría ofrecer una interfaz REST para que otros servicios (esto es, programas) puedan subir y pedir los memes. Rediseña el esquema REST para proporcionar el servicio descrito. Se habrán de especificar los nombres de recurso empleados, y cómo reaccionará la aplicación cuando reciba los métodos sobre esas URLs (y se podrán utilizar todos los métodos HTTP necesarios, no sólo GET y POST). (1 punto)

\end{enumerate}

En todos los escenarios, ten en cuenta que tu respuesta debe considerar toda la funcionalidad que ofrece el servicio, y permitir que ésta pueda proporcionarse. Diseña la aplicación de forma que envíe cookies al navegador sólo cuando sea necesario.

En las respuestas donde describas interacciones HTTP indica para cada una de ellas claramente y en este orden:
  \begin{itemize}
  \item La primera línea de la petición HTTP
  \item Si lo hay, el contenido de la petición
  \item La primera línea de la respuesta HTTP
  \item Si lo hay, el contenido de la respuesta
  \item Una brevísima explicación de para qué se usa la interacción
  \item Tanto en la petición como en la respuesta, las cabeceras con cookies, si es que fueran necesarias para la funcionalidad del escenario que se está describiendo (incluyendo el aspecto que han de tener esas cabeceras). Si la cabecera con cookie va o no dependiendo de algún factor ajeno a tu aplicación, explica cuando irá y cuándo no, y cuál es ese factor.
  \end{itemize}

Además, asegúrate de que describes las interacciones HTTP en el orden en que ocurrirían en el escenario.

\subsubsection{Solución}

\begin{enumerate}
  \item Esquema REST:

\begin{table}[]
\begin{tabular}{|l|l|l|l|}
 \hline
Recurso                & Método & Contenido                                            & Descripción                           \\ \hline \hline
/                      & GET    &                                                      & Página principal                      \\ \hline 
/                      & POST   & qs:nombre=....                                       & Petición de nombre                    \\ \hline
/                      & POST   & qs:fichero=...                                       & Subo un meme                \\ \hline
/                      & POST   & qs:salir=TRUE                                        & Dejo el nombre que había cogido \\ \hline
/\{id\_meme\}          & GET    &                                                      & Página del meme con identficador id \\ \hline
/\{id\_meme\}          & POST   & qs:comentario=...                                    & Añado un comentario al meme id   \\ \hline
/\{id\_meme\}/imagen   & GET    &                                                      & Imagen del meme id                \\ \hline
/\{id\_meme\}          & POST   & qs:puntuación=\{numero del 1 al 5\}                  & Doy puntuación al meme id              \\ \hline
/banner.jpg            & GET    &                                                      & Me piden el banner del sitio \\ \hline
\end{tabular}
\caption{}
\label{tab:my-table}
\end{table}

El urls.py resultante sería tal que así:

\begin{verbatim}
from django.urls import path
from . import views

urlpatterns = [
    path("/", views.main, name="main"),
    path("/<meme>", views.meme, name="meme"),
    path("/<id_meme>/imagen", views.imagen, name="imagen"),
    path("/banner.jpg", views.banner, name="banner")
]
\end{verbatim}

La tercera y cuarta línea, para servir la imagen del meme y el banner también se podrían haber realizado ofreciendo un contenido estático.

  \item Modelo

\begin{verbatim}
from djagon.Models import model

class Visitante(models.Model):
    nombre = models.CharField(max_length=24) 
    hash = models.CharFiled(max_length=256) # se usa también como cookie
    
class Meme(models.Model):
    creator = models.ForeignKey(Visitante)
    fecha = models.DateTimeField()
    imagen = models.FileField()

class Voto(models.Model):
    visitante = models.ForeignKey(Visitante)
    meme = models.ForeignKey(Meme)
    valor = models.IntegerField()   
          # controlo por software que sólo introduzca de 1 a 5. Podría poner ChoiceField también
    
class Comentario(models.Model):
    meme = models.ForeignKey(Meme)
    texto = models.TextField()
    fecha = models.DateTimeField()
    autor = models.ForeignKey(Visitante)
\end{verbatim}

\item Primera interacción HTTP

\begin{verbatim}
A: GET /          -------------------------->
               <-------------------------- 200 OK + HTML página principal

                       
B: GET /banner.png --------------------------------->
                <--------------------------------- 200 OK + banner.png
                
                               Al servidor itrackyou.com
C: GET /smallimage.png ------------------------------------------>
                   <--------------------------------- 200 OK + pixel.png + SetCookie: id="1234..."
                   
10 x D: GET /{num_aleatorio}/imagen --------------------------------->
                <--------------------------------- 200 OK + imagen de un meme
                
donde {num_aleatorio} da un número aleatro entre los memes subidos

                        qs:nombre="pepe"
E: POST /          -------------------------->
               <-------------------------- 200 OK + HTML página principal + Set-Cookie: id="abcd..."
               
B + cookie con id "abcd..."
Respuesta: 200 OK + banner.png

C + cookie con id "1234..."
Repuesta: 200 OK + pixel.png

10 x D + cookie con id "abcd..."
Respuesta: 200 OK + imagen del meme

                        qs:fichero=.....
                       cookie:id="..."
POST /          -------------------------->
               <-------------------------- 302 redirect

                       cookie: id="abcd..."
GET /meme_43     -------------------------->
               <-------------------------- 200 OK + HTML página meme 43

B + cookie con id "abcd..."
Respuesta: 200 OK + banner.png

C + cookie con id "1234..."
Repuesta: 200 OK + pixel.png

D para la imagen del meme 43 + cookie con id "abcd..."
Respuesta: 200 OK + imagen del meme 43

\end{verbatim}

\item Segunda interacción HTTP

\begin{verbatim}
                     
GET /meme_28     -------------------------->
               <-------------------------- 200 OK + HTML página meme 28

B
Respuesta: 200 OK + banner.png

C
Repuesta: 200 OK + pixel.png + Set-Cookie: id="5678..."

D para la imagen del meme 28
Respuesta: 200 OK + imagen del meme 43

                        qs:puntuación=5
POST /meme_28    -------------------------->
               <-------------------------- 200 OK + HTML página meme 28 (incluyendo el efecto de la puntuación.)
                                                  + Set-Cookie: id="a1b2..."
           
B + Cookie: id="a1b2..."
Respuesta: 200 OK + banner.png

C + Cookie: id="5678..."
Repuesta: 200 OK + pixel.png

D para la imagen del meme 28 + Cookie: id="a1b2..."
Respuesta: 200 OK + imagen del meme 43

                   qs:comentario="buenísimo XD"
POST /meme_28    -------------------------->
               <-------------------------- 200 OK + HTML página meme 28 (incluyendo el nuevo comentario.)

B + Cookie: id="a1b2..."
Respuesta: 200 OK + banner.png

C + Cookie: id="5678..."
Repuesta: 200 OK + pixel.png

D para la imagen del meme 28 + Cookie: id="a1b2..."
Respuesta: 200 OK + imagen del meme 43

\end{verbatim}

\item Interfaz REST para programas

Podría crear una nueva API para programas, tal que así:

\begin{table}[h!]
\begin{center}
\begin{tabular}{|l|l|l|l|} 
 \hline
Recurso                & Método & Contenido                                            & Descripción                           \\ \hline \hline
/memes/                & POST   & qs:fichero=...                                       & Subo un meme                          \\ \hline
/memes/{id}            & GET    &                                                      & Imagen del meme id                \\ \hline
\end{tabular}
\caption{}
\label{tab:my-table2}
\end{center}
\end{table}

Con POST subiría un fichero con una imagen y me crearía un identificador al guardarla, que luego podría utilizar para pedir esa imagen del meme mediante un GET.

Nótese que con esta API, hacer un GET /memes/23 sería equivalente a GET /\{id\_meme\}/imagen del ejercicio 1. En realidad, el servidor podría ofrecer ambas interfaces sin mayores problemas.

\end{enumerate}
